% 例 1（等腰存在性）：A(0,3), B(4,0)，C 在 x 轴；4 个解 C1=(7/8,0), C2=(-4,0), C3=(9,0), C4=(-1,0)
\documentclass[tikz,border=4pt]{standalone}
\usepackage{ctex}
\usepackage{amsmath}
\usetikzlibrary{calc}
\begin{document}
\begin{tikzpicture}[thick, scale=0.55]
  \draw[->, gray] (-6,0) -- (10.5,0) node[right, font=\footnotesize] {$x$};
  \draw[->, gray] (0,-1) -- (0,5) node[above, font=\footnotesize] {$y$};
  \foreach \i in {-4,-1,4,9} \node[font=\tiny, below, gray] at (\i,0) {$\i$};
  \node[font=\footnotesize, below left, gray] at (0,0) {$O$};
  \coordinate (A) at (0,3);
  \coordinate (B) at (4,0);
  \filldraw (A) circle (2pt) node[above left, font=\footnotesize] {$A(0,3)$};
  \filldraw (B) circle (2pt) node[below, font=\footnotesize] {$B(4,0)$};
  \draw[thick] (A) -- (B);
  % 4 个 C 点
  \foreach \x/\lbl in {0.875/$C_1$, -4/$C_2$, 9/$C_3$, -1/$C_4$} {
    \filldraw[red] (\x,0) circle (2pt);
    \node[red, font=\footnotesize, below] at (\x,-0.1) {\lbl};
  }
  % 几条解的连线（淡）
  \draw[red, dashed, very thin] (A) -- (0.875,0) -- (B);
  \draw[red, dashed, very thin] (A) -- (-4,0) -- (B);
  \draw[red, dashed, very thin] (A) -- (9,0) -- (B);
  \draw[red, dashed, very thin] (A) -- (-1,0) -- (B);
  \node[font=\footnotesize] at (3, 4.5) {等腰 $\triangle ABC$ 共 $4$ 解};
\end{tikzpicture}
\end{document}
